% !TEX root = ../main.tex
\chapter{極限・引き戻し：制約を満たす合流点}
\label{chap:limits-pullbacks}
\BookIndex{きょくげん}{極限}
\BookIndex{pullback}{Pullback}

\begin{chapterabstract}
極限と引き戻しを、ソフトウェア上の「整合条件を満たす合流点」として導入します。
二つのデータが同じキー、同じバージョン、同じ権限範囲を指しているかを確認しながら合成する場面で具体化します。
Lean 4 では、まず構造体と述語を使って「整合するペア」を表します。
このモデルを、DB join、複数 API レスポンスの統合、設定合成、DB スキーマ変更の保存性へ適用します。
\end{chapterabstract}

\begin{learninggoals}
\item 極限を「複数の条件を同時に満たす対象」として説明できます。
\item Pullbackを「共通の比較先で一致するペア」として読めます。
\item DB joinやAPI統合を、単なるペアではなく整合条件付きデータとしてモデル化できます。
\item Lean 4で、構造体と述語を使って整合するペアを表せます。
\item 変換後にも整合性が保存される、という性質を小さな定理として書けます。
\end{learninggoals}

\section{複数の制約を同時に満たす}

ソフトウェアでは、二つの値を単に並べるだけなら簡単です。
たとえば、プロフィール情報と注文情報をペアにする、設定ファイルAと設定ファイルBをまとめる、ユーザーAPIと権限APIのレスポンスを同時に扱う、といった処理です。
この程度であれば、積やレコードで十分に見えます。

しかし、実務で必要になるのは、通常は「任意のペア」ではありません。
同じユーザーを指すプロフィールと注文、同じ環境に属する設定の組、同じリクエストIDに対応する複数APIの結果が必要です。
つまり、二つの値をまとめる前に、それらが同じものについて語っているかを確認する必要があります。

図\ref{fig:ch23-pullback-join}では、プロフィールと注文をそれぞれ \texttt{userId} へ写し、値が一致する組だけを合流点に置きます。

\FigurePlaceholder{fig-ch23-pullback-join.png}{0.90}{プロフィールと注文の引き戻し}{fig:ch23-pullback-join}

この章の目的は、圏論の極限を網羅的に学ぶことではありません。
実務でよく現れる「整合条件付きの結合」を、Leanで扱える小さなモデルに落とすことが目的です。

\section{ソフトウェア上の問題：ペアにすればよいとは限らない}

次のような状況を考えます。

注文 DB の \texttt{Order} とプロフィール API の \texttt{Profile} を一画面で使うなら、
\texttt{Order.userId} と \texttt{Profile.userId} が一致する条件を組に含める必要があります。

型だけを見ると、次のようなペアも作れます。

\begin{quote}
\texttt{(profileOfUser10, orderOfUser20)}
\end{quote}

しかし、これは業務上誤った組み合わせです。
型としては \texttt{Profile \ensuremath{\times} Order} に入りますが、表示してよい注文詳細ではありません。

この問題は、DB joinだけに限りません。

\begin{center}
\small
\setlength{\tabcolsep}{3pt}
\begin{tabular}{@{}p{0.17\linewidth}p{0.42\linewidth}p{0.35\linewidth}@{}}
\toprule
場面 & 単なるペアでは足りない理由 & 必要な整合条件 \\
\midrule
DB join & 違う行同士を結合できてしまう & 主キー・外部キーの一致 \\
複数API統合 & 別リクエストの結果を混ぜる恐れがある & ユーザーID、リクエストIDの一致 \\
設定合成 & 環境が違う設定を混ぜる恐れがある & 環境名、バージョンの一致 \\
権限つき表示 & データと権限が別ユーザーを指す恐れがある & 対象ID、権限スコープの一致 \\
\bottomrule
\end{tabular}
\end{center}

ここでは、「結合する」ことそのものではなく、「何によって整合していると言えるのか」を明示することが重要です。
Pullbackは、この「共通の比較先で一致する」という構図を抽象化した言葉です。

\section{極限を「条件を同時に満たすもの」として読む}

圏論でいう極限は、積、等化子、Pullbackなど、多くの構成を統一的に扱う広い概念です。
本章では、まず次の直感を押さえます。

極限は、図式に対する錐のうち、ほかの任意の錐から一意な媒介射を受け取る終対象です。
本章ではこの定義を、「複数の条件を同時に満たし、同じ条件を満たす任意のデータが一意に経由する合流点」と読みます。
条件を満たすだけでは極限とは言えず、媒介射の存在、可換性、一意性までが必要です。

たとえば、積は二つの情報を同時に持つ対象です。

\[
A \times B
\]

これは、\texttt{A} と \texttt{B} の両方を持つという条件を満たします。
しかし、\texttt{A} と \texttt{B} の間に追加の整合条件はありません。
したがって、任意の \texttt{A} と任意の \texttt{B} を組み合わせられます。

型と関数の圏における pullback では、もう一つ条件が加わります。
\texttt{A} と \texttt{B} がそれぞれ共通の比較先 \texttt{C} に写るとき、その比較結果が一致するものだけを集めます。

\[
A \times_C B \cong \{(a,b) \mid f(a) = g(b)\}
\]

ここで \(\cong\) としたのは、右辺が型の要素を使った具体的な表示だからです。
一般の圏では要素の集合として定義せず、射影と普遍性で pullback を特徴付けます。
\texttt{C} は、ユーザー ID や環境名など、整合性を判定する共通の型です。

\section{Pullbackを「整合するペア」として読む}

Pullbackの基本形は、次のような図式で表されます。

\[
A \xrightarrow{f} C \xleftarrow{g} B
\]

このとき、Pullbackは「\texttt{A} の値と \texttt{B} の値のペアで、\texttt{C} に写した結果が一致するもの」です。

写像 \(f : A \to C\) と \(g : B \to C\) があるとします。
このとき Pullback は、直感的には次のようなデータです。

\[
\{(a,b) \mid f(a) = g(b)\}
\]

ソフトウェア上は、\texttt{A} と \texttt{B} の値を同時に持ち、さらに「共通キーが一致している」という証拠を持つ構造体として読めます。

この定義には、二つの要点があります。

第一に、Pullbackは単なるペアではなく、ペアに加えて整合条件の証拠を持ちます。

第二に、整合条件は外から後付けするコメントではありません。
型や構造の一部として持たせます。
Leanでは、この「証拠を持つ」という表現を直接記述できます。

DB join を pullback 的に読むなら、片側ずつキーを取り出す次の関数を考えます。
\begin{itemize}
\item \texttt{orderToUserId : Order \ensuremath{\to} UserId}
\item \texttt{profileToUserId : Profile \ensuremath{\to} UserId}
\end{itemize}
この具体例での pullback は、両関数の値が等しい注文とプロフィールの組です。

図\ref{fig:ch23-universal-property}では、整合する二つの写像を持つ任意の \texttt{X} から pullback への媒介写像があり、二つの射影を保つ写像はそれ一つに決まることを示します。

\FigurePlaceholder{fig-ch23-universal-property.png}{0.90}{pullback の普遍性}{fig:ch23-universal-property}

\subsection{普遍性を怖がらずに読む}

標準的なPullbackの普遍性は、可換条件を満たす媒介射の存在と一意性からなります~\cite{leinster-basic-category-theory,riehl-category-theory-in-context}。

Pullbackには、普遍性という条件があります。
実務上は、次のように読めます。

何らかのデータ \texttt{X} から \texttt{Profile} と \texttt{Order} を取り出せ、そのユーザー ID が常に一致するなら、
\texttt{X} から整合する組への写像を作れます。
普遍性は、射影を保つその写像が一意であることまで要求します。

この見方は、設計上も役立ちます。
複数の入力を統合する関数を書くとき、出力型が単なるペアなら、整合条件は呼び出し側の暗黙の約束になります。
Pullback的な型を使えば、その約束を出力型に含められます。

\section{Lean 4 で表す：まず構造体と述語で十分}

本章では、Mathlibの \texttt{CategoryTheory} API を使って一般のPullbackを形式化しません。
まず、実務で使える小さな定義から始めます。

Lean では、通常のフィールドに二つのデータを置き、追加フィールド
\texttt{ProfileOrder.consistent} にユーザー ID が等しい証明を置きます。

\begin{listing}[htbp]
\begin{minted}{lean4}
namespace Chapter23

abbrev UserId := Nat

structure Profile where
  userId : UserId
  displayName : String
deriving Repr, DecidableEq

structure Order where
  orderId : Nat
  userId : UserId
  total : Nat
deriving Repr, DecidableEq

structure ProfileOrder where
  profile : Profile
  order : Order
  consistent : profile.userId = order.userId
\end{minted}
\caption{\texttt{ProfileOrder}}
\label{lst:ch23-profile-order}
\end{listing}

\texttt{ProfileOrder} は、\texttt{Profile} と \texttt{Order} を持つだけではありません。
\texttt{consistent} というフィールドに、\texttt{profile.userId = order.userId} の証明を持っています。

この形にすると、以後の関数は「整合しているはずです」と仮定する必要がありません。
値そのものが整合性の証拠を含むためです。

\begin{listing}[htbp]
\begin{minted}{lean4}
def joinProfileOrder (p : Profile) (o : Order)
    (h : p.userId = o.userId) : ProfileOrder :=
  { profile := p, order := o, consistent := h }

theorem joined_profile_id_eq_order_id
    (j : ProfileOrder) : j.profile.userId = j.order.userId :=
  j.consistent

def orderToUserId (o : Order) : UserId := o.userId

def profileToUserId (p : Profile) : UserId := p.userId

theorem pullback_square_commutes (j : ProfileOrder) :
    profileToUserId j.profile = orderToUserId j.order :=
  j.consistent
\end{minted}
\caption{\texttt{pullback\_square\_commutes}}
\label{lst:ch23-pullback-square}
\end{listing}

\texttt{pullback\_square\_commutes} は、Pullbackで現れる可換四角形を、Leanの等式として書いたものです。
整合するペアからプロフィール側に射影してユーザーIDを取り出しても、注文側に射影してユーザーIDを取り出しても、同じ値になることを述べています。

\subsection{普遍性の最小モデル}

Pullbackの普遍性も、小さな関数として表せます。
任意の型 \texttt{X} から \texttt{Profile} と \texttt{Order} を取り出せて、かつそれらの \texttt{userId} が常に一致するなら、\texttt{X} から \texttt{ProfileOrder} への関数を作れます。

\begin{listing}[htbp]
\begin{minted}{lean4}
def toProfileOrder {X : Type}
    (f : X -> Profile) (g : X -> Order)
    (h : forall x, (f x).userId = (g x).userId) :
    X -> ProfileOrder :=
  fun x =>
    { profile := f x, order := g x, consistent := h x }

theorem toProfileOrder_profile {X : Type}
    (f : X -> Profile) (g : X -> Order)
    (h : forall x, (f x).userId = (g x).userId)
    (x : X) :
    (toProfileOrder f g h x).profile = f x := rfl

theorem toProfileOrder_order {X : Type}
    (f : X -> Profile) (g : X -> Order)
    (h : forall x, (f x).userId = (g x).userId)
    (x : X) :
    (toProfileOrder f g h x).order = g x := rfl

theorem toProfileOrder_unique {X : Type}
    (f : X -> Profile) (g : X -> Order)
    (h : forall x, (f x).userId = (g x).userId)
    (k : X -> ProfileOrder)
    (hp : forall x, (k x).profile = f x)
    (ho : forall x, (k x).order = g x) :
    k = toProfileOrder f g h := by
  funext x
  cases hx : k x with
  | mk p o consistent =>
      have hp' : p = f x := by simpa [hx] using hp x
      have ho' : o = g x := by simpa [hx] using ho x
      subst p
      subst o
      rfl
\end{minted}
\caption{\texttt{toProfileOrder}}
\label{lst:ch23-universal-map}
\end{listing}

このコードでは、普遍性の三要素を分けて確認しています。
\texttt{toProfileOrder} が媒介写像の存在、二つの射影定理が可換条件、
\texttt{toProfileOrder\_unique} がその条件を満たす媒介写像の一意性を表します。
一意性では、第\ref{chap:function-extensionality}章の関数外延性を使い、任意の入力における\texttt{profile}と\texttt{order}が一致することから関数全体の一致を示します。

ただし、これは\texttt{Type}と関数の世界における依存ペアのモデルです。
Mathlibの一般の圏における\texttt{IsLimit}やPullback APIを形式化したものではありません。
一般圏の定義との対応は、存在、可換性、一意性という証明義務の形に限られます。

\section{整合性を保存する変換}

整合するペアを作った後に、片側のデータを変換することがあります。
たとえば、注文金額への手数料の追加、注文表示用の金額の正規化、注文のメタ情報だけの更新などです。

このとき、変換が \texttt{userId} を変えないなら、整合性は保存されます。

注文の \texttt{total} だけを変更して \texttt{userId} を保存するなら、変換前の整合性の証明を変換後にも使えます。

\begin{listing}[htbp]
\begin{minted}{lean4}
def upgradeOrder (o : Order) : Order :=
  { o with total := o.total + 1 }

theorem upgradeOrder_preserves_userId (o : Order) :
    (upgradeOrder o).userId = o.userId := rfl

def liftJoinedOrder (j : ProfileOrder) : ProfileOrder :=
  { profile := j.profile,
    order := upgradeOrder j.order,
    consistent := by
      simpa [upgradeOrder] using j.consistent }

theorem liftJoinedOrder_preserves_consistency
    (j : ProfileOrder) :
    (liftJoinedOrder j).profile.userId =
      (liftJoinedOrder j).order.userId :=
  (liftJoinedOrder j).consistent
\end{minted}
\caption{\texttt{liftJoinedOrder}}
\label{lst:ch23-preserve-consistency}
\end{listing}

証明方針は単純です。
\texttt{upgradeOrder} は \texttt{total} だけを変え、\texttt{userId} を変えません。
したがって、元の \texttt{j.consistent} を、変換後の注文にも使える形へ書き換えます。

このような性質は、スキーマ移行や表示用変換でよく現れます。
変換関数がキーを保存することを先に証明しておくと、join結果や整合データ全体にも安全に持ち上げられます。

\section{実務に近いミニケース：注文、プロフィール、配送}

二つの情報源だけでなく、三つ以上の情報源を照合することもあります。
注文、プロフィール、配送情報をまとめる場合、少なくとも次のような条件が必要です。

\begin{itemize}
\item 注文とプロフィールの \texttt{userId} が一致します。
\item 配送情報と注文の \texttt{orderId} が一致します。
\item 配送情報と注文の \texttt{userId} も一致します。
\end{itemize}

ここでは、まず注文と配送の整合性だけを小さく切り出します。

\begin{listing}[htbp]
\begin{minted}{lean4}
structure Shipment where
  orderId : Nat
  userId : UserId
  delivered : Bool
deriving Repr, DecidableEq

def shipmentMatchesOrder (s : Shipment) (o : Order) : Prop :=
  s.userId = o.userId /\ s.orderId = o.orderId

structure OrderShipment where
  order : Order
  shipment : Shipment
  matched : shipmentMatchesOrder shipment order

theorem orderShipment_userId_eq (os : OrderShipment) :
    os.shipment.userId = os.order.userId :=
  os.matched.left

theorem orderShipment_orderId_eq (os : OrderShipment) :
    os.shipment.orderId = os.order.orderId :=
  os.matched.right
\end{minted}
\caption{\texttt{OrderShipment}}
\label{lst:ch23-order-shipment}
\end{listing}

\texttt{shipmentMatchesOrder} は、二つの条件を同時に持つ述語です。
Leanでは \verb|/\| が「かつ」を表します。
この構造により、\texttt{OrderShipment} の値を持っているだけで、ユーザーIDの一致と注文IDの一致の両方を取り出せます。

図\ref{fig:ch23-api-db-consistency}は、プロフィール、注文、配送情報を一つのレコードへ無条件に集めず、
ユーザー ID と注文 ID の一致をそれぞれ証明義務にする構成を示します。

\FigurePlaceholder{fig-ch23-api-db-consistency.png}{0.92}{複数データ源の整合性}{fig:ch23-api-db-consistency}

実務では、テナントID、リージョン、データのバージョン、権限スコープ、削除済みフラグなど、制約がさらに増えることがあります。
Pullbackの見方は、そうした制約を「何となく確認するもの」ではなく、「合流点の型に含めるもの」として整理するのに役立ちます。

\section{DB join、設定合成、複数APIレスポンスへの読み替え}

\subsection{DB join}

DB joinでは、二つのテーブルの行を共通キーで結び付けます。
SQLでは、\texttt{ON orders.user\_id = profiles.user\_id} のように条件を書きます。
この条件が、Pullback的な等式です。

Leanの小さなモデルでは、DBエンジンそのものは扱いません。
代わりに、「join結果として返してよい値とは何か」を型で表します。
この型により、スキーマ変更後にも join 結果が保存されるという仕様を明確に書けます。

\subsection{設定合成}

設定ファイルを複数の層から合成する場合にも、整合条件が必要です。

\begin{itemize}
\item base設定とoverride設定が同じ環境を指していますか。
\item 同じアプリケーション名ですか。
\item 対応するスキーマバージョンですか。
\end{itemize}

設定を単純にマージするだけなら、モノイド的に見えます。
しかし、マージしてよい組み合わせを制限するなら、Pullback的な整合条件が前面に出ます。

\subsection{複数APIレスポンス}

複数APIから情報を集める場合、レスポンスが同じユーザー、同じリクエスト、同じ時点を指しているかを確認する必要があります。

このとき、各レスポンスから共通の比較先へ写す関数を考えます。

\begin{itemize}
\item \texttt{profileResponseToUserId}
\item \texttt{permissionResponseToUserId}
\item \texttt{billingResponseToUserId}
\end{itemize}

それらの結果が一致しているときだけ、統合レスポンスを作ります。
この統合レスポンスは、単なるレコードではなく、整合性の証拠を持つレコードです。

\section{よくある誤解}

Pullbackは「joinの別名」ではありません。
joinは実装やクエリの操作であり、Pullbackは「共通の比較先で一致するものだけを集める」という構造の見方です。

また、本章のLeanコードは、実DBの完全な意味論を証明しているわけではありません。
インデックス、NULL、トランザクション、分散整合性、権限チェック、パフォーマンスなどは別の問題です。
本章で証明しているのは、モデル化した範囲で、整合条件が型と命題として明示されているという点です。

もう一つの誤解は、「整合条件はテストで十分である」という考え方です。
テストは必要ですが、テストだけでは代表例しか確認できません。
Leanで整合条件を型に含めると、その型の値を作る時点で証拠が必要になります。
これにより、整合していない値を後続処理へ渡す経路を減らせます。

\section{本章のまとめ}

極限は、本章では「複数の条件を同時に満たす合流点」として読みました。

Pullbackは、\(A \to C \leftarrow B\) という形に対し、\texttt{A} と \texttt{B} の値が \texttt{C} 上で一致するペアとして読めます。

ソフトウェアでは、PullbackはDB join、複数APIレスポンスの統合、設定合成、データ照合の直感に対応します。

Leanでは、一般の圏論APIを使わなくても、構造体と述語で「整合するペア」を表せます。

変換が共通キーを保存するなら、整合するペア全体へ安全に持ち上げられます。

本章の証明は、実システム全体の完全検証ではなく、仕様の核を小さく明示するためのモデルです。
